% ============================================================
%  Übungsaufgaben Template (my_dhbw Stil, uebung_*.tex)
%  Header "Übungsblatt", grün eingefärbte <details>-Lösungen.
%  Renderer: pdflatex (zweimal)
% ============================================================
\documentclass[11pt, a4paper]{article}

\usepackage[ngerman]{babel}
\usepackage{fontspec}
\setmainfont[
  Path=/usr/share/fonts/TTF/,
  Extension=.ttf,
  BoldFont=DejaVuSans-Bold,
  ItalicFont=DejaVuSans-Oblique,
  BoldItalicFont=DejaVuSans-BoldOblique
]{DejaVuSans}
\setsansfont[
  Path=/usr/share/fonts/TTF/,
  Extension=.ttf,
  BoldFont=DejaVuSans-Bold,
  ItalicFont=DejaVuSans-Oblique,
  BoldItalicFont=DejaVuSans-BoldOblique
]{DejaVuSans}
\setmonofont[Path=/usr/share/fonts/TTF/,Extension=.ttf]{DejaVuSansMono}
\usepackage[top=2cm, bottom=2cm, left=2.4cm, right=2.4cm, footskip=1cm]{geometry}
\usepackage{amsmath, amssymb}
\usepackage{xcolor}
\usepackage{enumitem}
\usepackage{fancyhdr}
\usepackage{lastpage}
\usepackage{tabularx}
\usepackage{booktabs}
\usepackage{microtype}
\usepackage{parskip}

\usepackage{array}
\usepackage{longtable}
\usepackage{ltablex}
\keepXColumns
\usepackage{colortbl}
\usepackage[most,breakable]{tcolorbox}
\usepackage{listings}
\usepackage[normalem]{ulem}
\usepackage{pdflscape}
\usepackage{titlesec}
\usepackage[colorlinks=true, linkcolor=accent, urlcolor=accent]{hyperref}

\renewcommand{\familydefault}{\sfdefault}

% ── Theme ─────────────────────────────────────────────────────
\definecolor{accent}{RGB}{0, 60, 113}
\definecolor{primaryblue}{RGB}{0, 60, 113}
\definecolor{loesung}{RGB}{0, 100, 0}
\definecolor{codebg}{RGB}{245, 245, 245}
\definecolor{figurebg}{RGB}{232, 241, 255}
\definecolor{tablehintbg}{RGB}{240, 255, 240}
\definecolor{solutionbg}{RGB}{240, 252, 240}
\definecolor{rulegray}{RGB}{180, 180, 180}
\definecolor{tableheadbg}{RGB}{0, 60, 113}

% ── Header/Footer ─────────────────────────────────────────────
\pagestyle{fancy}
\fancyhf{}
\renewcommand{\headrulewidth}{0.4pt}
\fancyhead[L]{\footnotesize\color{accent}\nichttitle}
\fancyhead[R]{\footnotesize\color{accent}\nichtsubtitle}
\fancyfoot[R]{\footnotesize\color{accent}Blatt \thepage\,/\,\pageref{LastPage}}

% ── Typografie ────────────────────────────────────────────────
\titleformat{\section}
    {\color{accent}\large\bfseries}{}{0pt}{}
    [\vspace{-4pt}\color{accent}\rule{\linewidth}{1.5pt}]
\titleformat{\subsection}
    {\normalsize\bfseries\color{accent!85!black}}{}{0pt}{}{}
\titleformat{\subsubsection}
    {\small\bfseries\color{accent!70!black}}{}{0pt}{}{}
\titlespacing*{\section}{0pt}{10pt}{3pt}
\titlespacing*{\subsection}{0pt}{6pt}{2pt}
\titlespacing*{\subsubsection}{0pt}{4pt}{1pt}

\setlength{\parindent}{0pt}
\setlength{\parskip}{6pt}
\renewcommand{\arraystretch}{1.2}
\setlist[itemize]{noitemsep, topsep=3pt, parsep=0pt, leftmargin=1.4em}
\setlist[enumerate]{noitemsep, topsep=3pt, parsep=0pt, leftmargin=1.6em}

% ── Boxen: solutionbox = Lösungsbox in grün ──────────────────
\newtcolorbox{figurebox}[1]{
  colback=figurebg, colframe=accent!50,
  title={Abbildung: #1}, fonttitle=\small\bfseries,
  breakable, before skip=6pt, after skip=6pt
}
\newtcolorbox{tablehintbox}[1]{
  colback=tablehintbg, colframe=green!50!black!50,
  title={Tabelle: #1}, fonttitle=\small\bfseries,
  breakable, before skip=6pt, after skip=6pt
}
\newtcolorbox{solutionbox}[1]{
  enhanced,
  colback=solutionbg, colframe=loesung,
  title={#1}, fonttitle=\small\bfseries,
  coltitle=white, colbacktitle=loesung,
  breakable, before skip=8pt, after skip=8pt
}

\lstset{
  basicstyle=\ttfamily\small,
  backgroundcolor=\color{codebg},
  breaklines=true,
  frame=single, framerule=0pt,
  xleftmargin=4pt, xrightmargin=4pt,
  aboveskip=6pt, belowskip=6pt,
  keepspaces=true
}

\providecommand{\nichttitle}{Übungsblatt}
\providecommand{\nichtsubtitle}{}

\renewcommand{\nichttitle}{Digitale Betriebswirtschaftslehre}
\renewcommand{\nichtsubtitle}{Übungsblatt 1}


\begin{document}

\begin{center}
{\Large\bfseries\color{accent}\nichttitle}
\ifx\nichtsubtitle\empty\else
  \\[2pt]{\normalsize\color{accent!80!black}\nichtsubtitle}
\fi
\\[2pt]
\color{accent}\rule{0.4\linewidth}{0.8pt}\color{black}
\end{center}
\vspace{4pt}

% ── BODY ──────────────────────────────────────────────────────

\section{Übungsblatt 1 — Digitale Betriebswirtschaftslehre — Gesamtzettel}


\noindent\textcolor{rulegray}{\rule{\linewidth}{0.4pt}}


\paragraph{Aufgabe 1 — Produktionsfaktoren und Daten \textit{(15 Punkte)}}\mbox{}\\[2pt]

Die fiktive Firma \textit{SchraubFix GmbH} fertigt Spezialschrauben. Sie überlegt, ihre klassische Produktion durch eine datengetriebene Steuerung (Sensoren an Maschinen, Predictive Maintenance, Bedarfsprognose aus Verkaufsdaten) zu ergänzen.


a) Nennen Sie die drei klassischen Produktionsfaktoren nach Gutenberg und ordnen Sie folgende Ressourcen der SchraubFix GmbH zu: \textit{Drahtrollen aus Stahl, CNC-Drehmaschine, Lohnarbeiter, ERP-Software-Lizenz, Werkstattmeister}. \textit{(4 Punkte)}


b) Erklären Sie, warum \textbf{Daten} als moderner Produktionsfaktor gelten. Skizzieren Sie dazu den Wandel des klassischen Wertschöpfungsmodells (Input → Produktion → Output) hin zum datengetriebenen Modell. \textit{(5 Punkte)}


c) Begründen Sie, warum eine reine Klassifikation der Sensor-Daten als „Betriebsmittel" zu kurz greift. Welche Konsequenz hat die Anerkennung von Daten als eigener Produktionsfaktor für die Investitionsplanung der SchraubFix GmbH? \textit{(6 Punkte)}


\textbf{Lösung:}


a)

\begin{itemize}
  \item \textbf{Arbeit}: Lohnarbeiter, Werkstattmeister
  \item \textbf{Betriebsmittel (Kapital)}: CNC-Drehmaschine, ERP-Software-Lizenz
  \item \textbf{Werkstoffe}: Drahtrollen aus Stahl
\end{itemize}

b) Daten tragen über die Kette \textbf{Sammlung → Verarbeitung → Entscheidung} zur Wertschöpfung bei. Sie ermöglichen bessere Prognosen, effizientere Prozesse und neue Geschäftsmodelle. Wandel: das klassische Modell \textbf{Input → Produktion → Output} wird ergänzt/abgelöst durch \textbf{Daten → Analyse → Entscheidung → Wertschöpfung}. Daten sind nicht mehr nur Begleiterscheinung der Produktion, sondern eigenständige Quelle des Outputs.


c) Daten sind kein Betriebsmittel im klassischen Sinn, weil sie:

\begin{itemize}
  \item nicht durch Nutzung verbraucht werden (im Gegensatz zur Maschine, die abnutzt),
  \item beliebig duplizierbar sind,
  \item ihren Wert erst durch Verarbeitung und Entscheidung entfalten — nicht durch physische Transformation.
\end{itemize}

\textbf{Konsequenz für die Investitionsplanung}: Investitionen müssen nicht nur in Maschinen (Betriebsmittel) und Material (Werkstoffe), sondern eigenständig in Datenerfassung (Sensorik), Datenverarbeitung (Analytics, ML-Infrastruktur) und Daten-Kompetenz (Personal) fließen. Diese Investitionen werden direkt erfolgswirksam, weil sie Prognose- und Entscheidungsqualität verbessern — also unmittelbar die Wertschöpfung erhöhen.



\noindent\textcolor{rulegray}{\rule{\linewidth}{0.4pt}}


\paragraph{Aufgabe 2 — Marktgleichgewicht und Nachfrageschock \textit{(20 Punkte)}}\mbox{}\\[2pt]

Auf einem Studi-Markt für gebrauchte Tablets gelten:


\begin{itemize}
  \item Nachfrage: $Q_D = 240 - 4P$
  \item Angebot: $Q_S = 40 + 6P$
\end{itemize}

a) Berechnen Sie Gleichgewichtspreis $P^*$ und -menge $Q^*$. \textit{(5 Punkte)}


b) Die Hochschule diktiert einen Maximalpreis von $P = 15$ €. Bestimmen Sie $Q_D$, $Q_S$ und Art/Höhe des Marktungleichgewichts. Welche Marktreaktion wäre ohne Preisregulierung zu erwarten? \textit{(6 Punkte)}


c) Eine TikTok-Welle treibt die Nachfrage auf $Q_D' = 320 - 4P$. Berechnen Sie das neue Gleichgewicht. Was empfehlen Sie einem Anbieter aus Managementsicht? \textit{(5 Punkte)}


d) Erläutern Sie kurz, wie ein Plattformanbieter (z. B. Refurbished-Marktplatz) eine solche Nachfrageverschiebung \textbf{früher} erkennen könnte als ein klassischer Händler. Welcher Wettbewerbsvorteil entsteht? \textit{(4 Punkte)}


\textbf{Lösung:}


a) Gleichsetzen:

\[
240 - 4P = 40 + 6P \Rightarrow 200 = 10P \Rightarrow P^* = 20\,€
\]

\[
Q^* = 240 - 4 \cdot 20 = 160
\]


b) Bei $P = 15$ €:

\begin{itemize}
  \item $Q_D = 240 - 4 \cdot 15 = 180$
  \item $Q_S = 40 + 6 \cdot 15 = 130$
  \item $Q_D > Q_S$ → \textbf{Nachfrageüberschuss = 50 Tablets} (Preis zu niedrig).
\end{itemize}

Ohne Preisregulierung würde der Preis steigen, bis $P^* = 20$ € erreicht ist und der Markt geräumt wird.


c) Neues Gleichgewicht:

\[
320 - 4P = 40 + 6P \Rightarrow 280 = 10P \Rightarrow P^{*\prime} = 28\,€
\]

\[
Q^{*\prime} = 320 - 4 \cdot 28 = 208
\]


Die Nachfragekurve verschiebt sich nach rechts → \textbf{höherer Preis, höhere Menge}. Managementempfehlung: Kapazitäten/Beschaffung ausbauen und Verkaufspreis anheben (Richtung 28 €), um die zusätzliche Zahlungsbereitschaft abzuschöpfen.


d) Eine Plattform sieht Verkaufszahlen, Suchanfragen, Klickraten und Watchlists nahezu in Echtzeit über alle Anbieter und Käufer hinweg. Sie erkennt einen Anstieg der Nachfrage \textbf{bevor} Lagerbestände bei einzelnen Händlern leerlaufen. \textbf{Wettbewerbsvorteil}: schnellere Preis- und Mengenanpassung (Dynamic Pricing, gezielte Beschaffung) → höhere Marge und vermiedene Out-of-Stock-Situationen.



\noindent\textcolor{rulegray}{\rule{\linewidth}{0.4pt}}


\paragraph{Aufgabe 3 — Anbieterwahl unter Sicherheit (Transfer) \textit{(18 Punkte)}}\mbox{}\\[2pt]

Ein Online-Shop erwartet pro Jahr Transaktionen $x$ und wählt zwischen drei Payment-Anbietern:



{\small
\begin{tabularx}{\linewidth}{XXX}
\hline
\rowcolor{tableheadbg}
{\color{white}\textbf{Anbieter}} & {\color{white}\textbf{Gebühr pro Transaktion}} & {\color{white}\textbf{Fixkosten/Jahr}} \\[2pt]
\hline
\endhead
\hline
\endfoot
A (Basic) & 0,50 € & 5.000 € \\
B (Pro) & 0,35 € & 20.000 € \\
C (Enterprise) & 0,25 € & 50.000 € \\
\end{tabularx}
}


a) Stellen Sie für jeden Anbieter die Gesamtkostenfunktion $K(x)$ auf. \textit{(3 Punkte)}


b) Bestimmen Sie die kritischen Transaktionsmengen, ab denen B günstiger als A bzw. C günstiger als B wird. \textit{(6 Punkte)}


c) Empfehlen Sie für die Transaktionsvolumina $x_1 = 80\,000$, $x_2 = 200\,000$, $x_3 = 350\,000$ jeweils den optimalen Anbieter mit Begründung. \textit{(6 Punkte)}


d) Ordnen Sie diese Entscheidung in der Entscheidungstheorie ein: Sicherheit, Risiko oder Unsicherheit? Konstitutiv oder laufend? \textit{(3 Punkte)}


\textbf{Lösung:}


a)

\begin{itemize}
  \item $K_A(x) = 5.000 + 0{,}50\,x$
  \item $K_B(x) = 20.000 + 0{,}35\,x$
  \item $K_C(x) = 50.000 + 0{,}25\,x$
\end{itemize}

b) \textbf{A vs. B}: $5.000 + 0{,}50\,x = 20.000 + 0{,}35\,x \Rightarrow 0{,}15\,x = 15.000 \Rightarrow x = 100.000$.

\textbf{B vs. C}: $20.000 + 0{,}35\,x = 50.000 + 0{,}25\,x \Rightarrow 0{,}10\,x = 30.000 \Rightarrow x = 300.000$.


Also: bis 100.000 Transaktionen ist A günstig, zwischen 100.000 und 300.000 B, ab 300.000 C.


c)

\begin{itemize}
  \item $x_1 = 80.000$: $K_A = 45.000$, $K_B = 48.000$, $K_C = 70.000$ → \textbf{Anbieter A} (unter 100.000-Schwelle).
  \item $x_2 = 200.000$: $K_A = 105.000$, $K_B = 90.000$, $K_C = 100.000$ → \textbf{Anbieter B} (zwischen den Schwellen).
  \item $x_3 = 350.000$: $K_A = 180.000$, $K_B = 142.500$, $K_C = 137.500$ → \textbf{Anbieter C} (über 300.000-Schwelle).
\end{itemize}

d) Alle Kostenparameter sind vertraglich fixiert und das Volumen ist (annahmegemäß) bekannt → \textbf{Entscheidung unter Sicherheit}. Da der Anbieter nicht täglich gewechselt wird und die Wahl Geschäftsmodell und IT-Architektur prägt, handelt es sich um eine \textbf{konstitutive (oder zumindest mittelfristige strukturelle) Entscheidung}, nicht um eine laufende.



\noindent\textcolor{rulegray}{\rule{\linewidth}{0.4pt}}


\paragraph{Aufgabe 4 — Risiko, Unsicherheit und Hurwicz \textit{(20 Punkte)}}\mbox{}\\[2pt]

Ein Start-up plant den Markteintritt mit drei Alternativen. Das Ergebnis (Gewinn in T€) hängt vom Marktumfeld ab:



{\small
\begin{tabularx}{\linewidth}{XXXX}
\hline
\rowcolor{tableheadbg}
{\color{white}\textbf{Alternative}} & {\color{white}\textbf{Boom}} & {\color{white}\textbf{Stabil}} & {\color{white}\textbf{Rezession}} \\[2pt]
\hline
\endhead
\hline
\endfoot
A — Aggressiv & 300 & 100 & −200 \\
B — Moderat & 180 & 120 & −20 \\
C — Defensiv & 80 & 70 & 40 \\
\end{tabularx}
}


a) Die Geschäftsführung schätzt die Wahrscheinlichkeiten auf $p_{\text{Boom}}=0{,}3$, $p_{\text{Stabil}}=0{,}5$, $p_{\text{Rez}}=0{,}2$. Wählen Sie die optimale Alternative nach dem \textbf{Erwartungswert-Kriterium}. \textit{(5 Punkte)}


b) Ein neues Vorstandsmitglied bezweifelt diese Wahrscheinlichkeiten. Welche Alternative wäre nach \textbf{Maximin (Wald)}, \textbf{Maximax} und \textbf{Laplace} zu wählen? \textit{(6 Punkte)}


c) Wenden Sie das \textbf{Hurwicz-Kriterium} mit Optimismusparameter $\lambda = 0{,}4$ an und bestimmen Sie die optimale Alternative. \textit{(5 Punkte)}


d) Ordnen Sie die Vorgehensweisen aus a) und b) den Informationsständen \textit{Sicherheit / Risiko / Unsicherheit} zu und begründen Sie kurz, warum für ein junges Start-up das \textbf{Maximin-Kriterium} problematisch sein kann. \textit{(4 Punkte)}


\textbf{Lösung:}


a) Erwartungswert $E = \sum p_i \cdot x_i$:

\begin{itemize}
  \item $E(A) = 0{,}3 \cdot 300 + 0{,}5 \cdot 100 + 0{,}2 \cdot (-200) = 90 + 50 - 40 = 100$
  \item $E(B) = 0{,}3 \cdot 180 + 0{,}5 \cdot 120 + 0{,}2 \cdot (-20) = 54 + 60 - 4 = 110$
  \item $E(C) = 0{,}3 \cdot 80 + 0{,}5 \cdot 70 + 0{,}2 \cdot 40 = 24 + 35 + 8 = 67$
\end{itemize}

→ \textbf{Alternative B} (höchster Erwartungswert 110 T€).


b)

\begin{itemize}
  \item \textbf{Maximin}: jeweils Minimum → A: −200, B: −20, C: 40. Maximum davon → \textbf{C}.
  \item \textbf{Maximax}: jeweils Maximum → A: 300, B: 180, C: 80. Maximum davon → \textbf{A}.
  \item \textbf{Laplace} (Mittelwert je Zeile): A: $(300+100-200)/3 = 66{,}67$; B: $(180+120-20)/3 = 93{,}33$; C: $(80+70+40)/3 = 63{,}33$. Höchster Wert → \textbf{B}.
\end{itemize}

c) Hurwicz: $H = \lambda \cdot \max + (1-\lambda) \cdot \min$ mit $\lambda = 0{,}4$:

\begin{itemize}
  \item $H(A) = 0{,}4 \cdot 300 + 0{,}6 \cdot (-200) = 120 - 120 = 0$
  \item $H(B) = 0{,}4 \cdot 180 + 0{,}6 \cdot (-20) = 72 - 12 = 60$
  \item $H(C) = 0{,}4 \cdot 80 + 0{,}6 \cdot 40 = 32 + 24 = 56$
\end{itemize}

→ \textbf{Alternative B} (Hurwicz-Wert 60).


d) \textbf{Zuordnung}: a) Erwartungswert nutzt bekannte Wahrscheinlichkeiten → \textbf{Risiko}. b) Maximin/Maximax/Laplace verzichten auf Wahrscheinlichkeiten → \textbf{Unsicherheit}.


\textbf{Problem von Maximin für Start-ups}: Maximin ist extrem risikoavers — es maximiert nur das schlechteste Ergebnis und führt hier zur defensiven Alternative C mit nur 40–80 T€ Gewinn. Ein Start-up muss aber typischerweise wachsen, um Investoren zu überzeugen und Skaleneffekte zu erschließen. Maximin würde Chancen systematisch ausschlagen und wäre mit der Wachstumsanforderung (Shareholder-Erwartung) inkompatibel.



\noindent\textcolor{rulegray}{\rule{\linewidth}{0.4pt}}


\paragraph{Aufgabe 5 — Zielbeziehungen, SMART und Principal-Agent \textit{(17 Punkte)}}\mbox{}\\[2pt]

Die \textit{MeinBike AG} (Online-Händler für E-Bikes) verfolgt unter anderem folgende Ziele:


\begin{itemize}
  \item (Z1) „Kundenservice verbessern" durch Aufbau einer 24/7-Hotline
  \item (Z2) Reduktion der stationären Filialen um 50 \% bis Ende 2027
  \item (Z3) Erhöhung der Lagerumschlagshäufigkeit (Just-in-Time-Beschaffung)
  \item (Z4) Reduktion der Lagerkosten um 20 \%
  \item (Z5) Steigerung des EBIT auf mindestens 12 Mio. €
\end{itemize}

a) Klassifizieren Sie Z1–Z5 nach \textbf{Sach- vs. Formalzielen}. \textit{(4 Punkte)}


b) Bestimmen Sie die Zielbeziehung (komplementär / konkurrierend / indifferent) für die Paare \textbf{(Z1, Z2)} und \textbf{(Z3, Z4)}. Begründen Sie. \textit{(4 Punkte)}


c) Formulieren Sie Z1 nach \textbf{SMART} um. Geben Sie für jeden Buchstaben einen konkreten Aspekt an. \textit{(5 Punkte)}


d) Der Bereichsleiter Logistik wird mit einem Bonus von 20 \% seines Gehalts belohnt, wenn Z4 (−20 \% Lagerkosten) erreicht wird — unabhängig von anderen Zielen. Identifizieren Sie das \textbf{Principal-Agent-Problem} und schlagen Sie eine \textbf{bessere Anreizgestaltung} vor. \textit{(4 Punkte)}


\textbf{Lösung:}


a)

\begin{itemize}
  \item \textbf{Sachziele} (konkretes Handeln in betrieblichen Funktionen): \textbf{Z1, Z2, Z3} (Leistungs-/Organisationsebene).
  \item \textbf{Formalziele} (übergeordnete Erfolgsziele): \textbf{Z4} (Kostenziel/Wirtschaftlichkeit) und \textbf{Z5} (Gewinn/EBIT).
\end{itemize}
\textit{Hinweis}: Z4 ist ein wertmäßiges Erfolgsziel (Aufwandsseite) und damit Formalziel; Z3 ist Sachziel im Logistik-Prozess.


b)

\begin{itemize}
  \item \textbf{(Z1, Z2)}: \textbf{konkurrierend}. Eine 24/7-Hotline soll Servicequalität erhöhen, der Abbau stationärer Filialen reduziert aber Service-Touchpoints — analog zum Kapitelbeispiel „Personalisierter Service ↔ Reduktion stationäres Service-Netz".
  \item \textbf{(Z3, Z4)}: \textbf{komplementär}. Höhere Lagerumschlagshäufigkeit (Just-in-Time) reduziert Lagerbestände und damit Lagerkosten — analog zum Kapitelbeispiel „Just-in-Time ↔ Reduktion Lagerkosten".
\end{itemize}

c) \textbf{SMART-Formulierung von Z1}:

\begin{itemize}
  \item \textbf{S} (spezifisch): Aufbau einer 24/7-Telefon- und Chat-Hotline für E-Bike-Kund:innen.
  \item \textbf{M} (messbar): Antwortzeit ≤ 60 Sekunden, Lösungsquote im Erstkontakt ≥ 80 \%, NPS ≥ 50.
  \item \textbf{A} (attraktiv): stärkt Kundenbindung und differenziert von preisaggressiven Wettbewerbern.
  \item \textbf{R} (realistisch): mit dem geplanten Service-Budget (z. B. 1,5 Mio. € p. a.) und 25 zusätzlichen Mitarbeitenden machbar.
  \item \textbf{T} (terminiert): vollständig live bis 31.12.2026.
\end{itemize}

d) \textbf{Principal-Agent-Problem}: Der Vorstand (Principal) kann nicht im Detail beobachten, \textit{wie} der Bereichsleiter (Agent) die Lagerkosten senkt. Da der Bonus nur an Z4 gekoppelt ist, hat der Agent einen Anreiz, Lagerbestände unter ein vernünftiges Sicherheitsniveau zu drücken — was Lieferfähigkeit und Umsatz (Z5!) gefährdet. Klassische \textbf{Informationsasymmetrie + einseitige Anreize} führen zu einer Suboptimierung auf Kosten des Gesamtunternehmens.


\textbf{Bessere Anreizgestaltung}: Mehrdimensionale KPIs in der Bonusformel — z. B. 50 \% Lagerkostensenkung, 30 \% Service-Level (Lieferquote ≥ 98 \%), 20 \% EBIT-Beitrag des Bereichs. Ergänzend: Monitoring der Lieferfähigkeit, Cap auf Bonusauszahlung bei Verfehlen einer Mindest-Servicegrenze. So wird das Eigeninteresse des Agenten an die Gesamtziele (Stakeholder + Shareholder) gekoppelt.



\noindent\textcolor{rulegray}{\rule{\linewidth}{0.4pt}}


\paragraph{Aufgabe 6 — Effizienz, Effektivität und Kennzahlen (Code-Analyse) \textit{(15 Punkte)}}\mbox{}\\[2pt]

Ein Studi-Team schreibt ein Pseudocode-Skript, um die Wirtschaftlichkeit eines Produktionslaufs auszuwerten:


\begin{lstlisting}
input_menge       = 50           # kg Draht
output_menge      = 4000         # Schrauben
draht_preis_pro_kg = 2.0         # €
verkaufspreis     = 0.02         # €/Schraube

produktivitaet  = input_menge / output_menge
aufwand         = input_menge * draht_preis_pro_kg
ertrag          = output_menge * verkaufspreis
wirtschaftlich  = aufwand / ertrag
gewinn          = aufwand - ertrag

print("Produktivitaet:", produktivitaet)
print("Wirtschaftlichkeit:", wirtschaftlich)
print("Gewinn:", gewinn)
\end{lstlisting}


a) Finden Sie \textbf{drei inhaltliche Fehler} in den Berechnungen und korrigieren Sie diese. \textit{(6 Punkte)}


b) Berechnen Sie mit korrigierten Formeln Produktivität, Wirtschaftlichkeit und Gewinn für die gegebenen Werte. \textit{(4 Punkte)}


c) Ein Konkurrent erreicht mit denselben 50 kg Draht 5.000 Schrauben (gleicher Verkaufs-/Einkaufspreis). Vergleichen Sie beide Betriebe nach \textbf{Effizienz} und \textbf{Effektivität}, wenn das Unternehmensziel „Verkauf von genau 4.000 Schrauben pro Lauf" lautet. Wer ist effizienter, wer effektiver? \textit{(5 Punkte)}


\textbf{Lösung:}


a) \textbf{Fehler 1}: \texttt{produktivitaet = input\_menge / output\_menge} ist invertiert. Korrekt: \texttt{Output / Input} → \texttt{output\_menge / input\_menge}.

\textbf{Fehler 2}: \texttt{wirtschaftlich = aufwand / ertrag} ist invertiert. Korrekt: \texttt{Ertrag / Aufwand}.

\textbf{Fehler 3}: \texttt{gewinn = aufwand - ertrag} hat falsches Vorzeichen. Korrekt: \texttt{gewinn = ertrag - aufwand}.


b) Mit korrigierten Formeln:

\begin{itemize}
  \item \textbf{Produktivität} = 4.000 / 50 = \textbf{80 Schrauben/kg}
  \item \textbf{Aufwand} = 50 · 2,00 = 100 €
  \item \textbf{Ertrag} = 4.000 · 0,02 = 80 €
  \item \textbf{Wirtschaftlichkeit} = 80 / 100 = \textbf{0,8} (< 1 → unwirtschaftlich)
  \item \textbf{Gewinn} = 80 − 100 = \textbf{−20 €} (Verlust)
\end{itemize}

c) \textbf{Konkurrent}: Produktivität = 5.000 / 50 = 100 Schrauben/kg; Ertrag = 5.000 · 0,02 = 100 €; Aufwand = 100 €; Wirtschaftlichkeit = 1,0.


\begin{itemize}
  \item \textbf{Effizienz} (Werden die Dinge richtig gemacht? — Verhältnis Leistung/Ressourceneinsatz): Der Konkurrent ist mit 100 Schrauben/kg vs. 80 Schrauben/kg \textbf{effizienter} — er nutzt denselben Input besser.
  \item \textbf{Effektivität} (Werden die richtigen Dinge gemacht? — Grad der Zielerreichung): Ziel sind genau 4.000 Schrauben. Das eigene Unternehmen produziert exakt 4.000 → \textbf{100 \% Zielerreichung, also effektiv}. Der Konkurrent produziert 5.000 → 25 \% Übererfüllung; gemessen am Sollziel ist er weniger effektiv (Überproduktion = falsches Ziel getroffen).
\end{itemize}

\textbf{Fazit}: Der Konkurrent arbeitet effizienter, das eigene Unternehmen effektiver. Dass Effizienz und Effektivität auseinanderfallen können, ist genau der Grund, warum beide Größen getrennt geführt werden müssen.





% ──────────────────────────────────────────────────────────────

\end{document}
